% AN ALIGNMENT CHANGES THE MODE BEFORE IT READS ITS SPEC.
%
% \halign to<dimen> and \valign to<dimen> are met in one mode and run
% their entries in another. The change happens FIRST, before the `to' or
% `spread' keyword is looked for, so the <dimen> after it is scanned in
% the mode the entries will run in -- not the mode the alignment was met
% in. \tracingcommands=2 says so, because it names the mode whenever the
% mode has changed since the line it last drew:
%
%   met in                  the spec is read in
%   vertical mode           internal vertical mode      \halign
%   internal vertical       internal vertical mode      \halign
%   display math            internal vertical mode      \halign
%   horizontal mode         restricted horizontal mode  \valign
%
% So `\halign to\the\hsize{...}' in vertical mode draws
%
%   {internal vertical mode: \the}
%
% and an \halign that replaces a display draws the same, whether or not
% the display had anything in it for the "Improper \halign inside $$'s"
% fault to throw away. The mode is already internal vertical by the time
% the fault is reported.
\catcode`\{=1 \catcode`\}=2 \catcode`\#=6 \catcode`\$=3 \catcode`\&=4
\tracingonline=1 \hsize=100pt \parindent=0pt \tabskip=0pt
\hbadness=10000 \vbadness=10000 \hfuzz=1000pt \vfuzz=1000pt
\message{[A \halign in vertical mode]}
\tracingcommands=2
\halign to\the\hsize{#&#\cr\cr}
\tracingcommands=0
\message{[B \valign in horizontal mode]}
\tracingcommands=2
\noindent A\valign to\the\hsize{#\cr\cr}B\par
\tracingcommands=0
\message{[C \halign in internal vertical mode]}
\tracingcommands=2
\setbox0=\vbox{\halign to\the\hsize{#&#\cr\cr}}
\tracingcommands=0
\message{[D \halign that replaces an empty display]}
\tracingcommands=2
$$\halign to\the\displaywidth{#&#\cr\cr}$$
\tracingcommands=0
\message{[E \halign after something in a display]}
\tracingcommands=2
$$x\halign to\the\displaywidth{#&#\cr\cr}$$
\tracingcommands=0
\message{[done]}
\end
